\begin{center}
\begin{tikzpicture}[x=1cm,y=1cm]
\draw[rounded corners=2pt,fill=blue!3,draw=blue!22] (0,0) rectangle (6.35,2.75);
\node[anchor=west,font=\small\bfseries,text=newsblue] at (.2,2.46) {中国 2026 年前九月财政指标同比};
\draw[gray!60] (2.45,.44) -- (2.45,2.14);
\foreach \name/\y/\val/\len in {财政收入/1.96/+5.1/1.28,税收收入/1.61/+6.1/1.53,财政支出/1.26/+1.0/.25,土地出让收入/.91/-27.9/-1.40} {
 \node[anchor=east,font=\scriptsize] at (2.3,\y) {\name};
}
\fill[newsblue!70] (2.45,1.87) rectangle (3.73,2.05);
\fill[newsblue!85] (2.45,1.52) rectangle (3.98,1.70);
\fill[newsblue!50] (2.45,1.17) rectangle (2.70,1.35);
\fill[orange!75] (1.05,.82) rectangle (2.45,1.00);
\node[anchor=west,font=\scriptsize] at (3.80,1.96) {$+5.1\%$};
\node[anchor=west,font=\scriptsize] at (4.05,1.61) {$+6.1\%$};
\node[anchor=west,font=\scriptsize] at (2.78,1.26) {$+1.0\%$};
\node[anchor=east,font=\scriptsize] at (.99,.91) {$-27.9\%$};
\node[anchor=west,font=\scriptsize,text=newsgray] at (.25,.25) {图示为各自同比变化；正负条使用不同示意比例，不宜互比长度。};
\end{tikzpicture}
\par\footnotesize\color{newsgray}图 2｜编辑部整理。数据来自财政部 10 月 23 日发布会，经路透社报道；土地出让收入与一般公共预算收入分属不同统计体系。\color{black}
\end{center}